\section*{Bab \ref{ch:b2:frontier-orbitals} --- Orbital Perbatasan dan Kereaktifan}

\begin{solution}{exo:b2:frontier-orbitals:1}
Etena: \omterm{def:b2:frontier-orbitals:frontier}{HOMO} $\alpha + \beta$, \omterm{def:b2:frontier-orbitals:frontier}{LUMO} $\alpha - \beta$. Anion alil (empat elektron): \omterm{def:b2:frontier-orbitals:frontier}{HOMO} $\alpha$
(tingkat nonikatan), \omterm{def:b2:frontier-orbitals:frontier}{LUMO} $\alpha - 1.414\beta$. Butadiena: \omterm{def:b2:frontier-orbitals:frontier}{HOMO} $\alpha + 0.618\beta$, \omterm{def:b2:frontier-orbitals:frontier}{LUMO} $\alpha -
0.618\beta$.
\end{solution}

\begin{solution}{exo:b2:frontier-orbitals:2}
Keras: \ce{OH-}, \ce{H+}, \ce{F-}. Lunak: \ce{I-}, \ce{RS-}, karbon iodometana.
\end{solution}

\begin{solution}{exo:b2:frontier-orbitals:3}
Butadiena (konformer s-cis dapat dicapai) dan siklopentadiena (terkunci
s-cis) bereaksi; penta-1,4-diena tidak terkonjugasi dan tidak bereaksi.
\omterm{def:b2:frontier-orbitals:diels-alder}{Diena} (2\textit{Z},4\textit{Z}) hanya dapat mencapai bentuk s-cis dengan
mendorong kedua gugus metil bagian dalamnya saling bertabrakan: \omterm{def:b2:frontier-orbitals:diels-alder}{diena} itu
bereaksi sangat buruk.
\end{solution}

\begin{solution}{exo:b2:frontier-orbitals:4}
Sikloheks-3-ena-1-karbonitril: cincin beranggota enam, ikatan rangkap antara
bekas C2 dan C3 \omterm{def:b2:frontier-orbitals:diels-alder}{diena}, gugus \ce{C#N} pada salah satu dari kedua karbon
cincin yang berasal dari \omterm{def:b2:frontier-orbitals:diels-alder}{dienofil}.
\end{solution}

\begin{solution}{exo:b2:frontier-orbitals:5}
$\Delta\varepsilon = \qty{8}{eV}$: perturbatif $2\beta^2/\Delta\varepsilon = \qty{0.25}{eV}$. Eksak: tingkat
bawah $-6 - \sqrt{16 + 1} = \qty{-10.123}{eV}$, keuntungan $2 \times 0.123 = \qty{0.246}{eV}$.
\end{solution}

\begin{solution}{exo:b2:frontier-orbitals:6}
Iodometana adalah \omterm{def:b2:frontier-orbitals:hard-soft}{elektrofil lunak}: \omterm{def:b2:frontier-orbitals:control}{kendali orbital}, serangan oleh atom
dengan koefisien \omterm{def:b2:frontier-orbitals:frontier}{HOMO} yang lebih besar, yaitu karbon. Karbokation bebas
adalah \omterm{def:b2:frontier-orbitals:hard-soft}{elektrofil keras}: \omterm{def:b2:frontier-orbitals:control}{kendali muatan}, serangan sebagian oleh nitrogen,
yang membawa muatan negatif lebih besar.
\end{solution}

\begin{solution}{exo:b2:frontier-orbitals:7}
C4 \omterm{def:b2:frontier-orbitals:diels-alder}{diena} (koefisien \omterm{def:b2:frontier-orbitals:frontier}{HOMO} terbesar) berikatan dengan karbon $\beta$ propenal
(koefisien \omterm{def:b2:frontier-orbitals:frontier}{LUMO} terbesar): 2-metoksisikloheks-3-ena-1-karbaldehida, dengan
kedua substituen pada karbon-karbon yang bertetangga (produk ``orto'').
\end{solution}

\begin{solution}{exo:b2:frontier-orbitals:8}
Kerangka bisiklo[2.2.1]heptena (norbornena) dengan gugus ester pada salah
satu karbon bekas \omterm{def:b2:frontier-orbitals:diels-alder}{dienofil}, mengarah ke ikatan \ce{C=C} yang baru, di bawah
cincin beranggota enam (endo), pada sisi yang berseberangan dengan jembatan
\ce{CH2}.
\end{solution}

\begin{solution}{exo:b2:frontier-orbitals:9}
Kedua gugus karbonil menurunkan \omterm{def:b2:frontier-orbitals:frontier}{LUMO} \omterm{def:b2:frontier-orbitals:diels-alder}{dienofil}: celahnya dengan \omterm{def:b2:frontier-orbitals:frontier}{HOMO}
siklopentadiena jauh lebih kecil daripada celah etena, stabilisasi
$\beta^2/\Delta\varepsilon$ keadaan transisinya lebih besar, dan penghalangnya
lebih rendah.
\end{solution}

\begin{solution}{exo:b2:frontier-orbitals:10}
Reaksinya stereospesifik: maleat menghasilkan adduk dengan kedua gugus ester
cis (endo,endo sebagai produk utama; adduk ini akiral, sebuah senyawa meso);
fumarat menghasilkan adduk trans, satu ester endo dan satu exo, yang
terbentuk sebagai sepasang enansiomer (rasemat).
\end{solution}

\begin{solution}{exo:b2:frontier-orbitals:11}
Dua ikatan $\pi$ menjadi dua ikatan $\sigma$ yang lebih kuat: $\Delta_r H^\circ < 0$. Dua molekul menjadi satu:
$\Delta_r S^\circ < 0$. $\Delta_r G^\circ = \Delta_r H^\circ - T\Delta_r S^\circ$ bertambah seiring $T$ dan menjadi
positif di atas $T_i = \Delta_r H^\circ/\Delta_r S^\circ$: pada suhu tinggi adduknya terpecah kembali
(retro-Diels--Alder).
\end{solution}

\begin{solution}{exo:b2:frontier-orbitals:12}
Dua orbital pasangan elektron bebas yang terisi dan bersebelahan membentuk
interaksi empat elektron, yang tolak-menolak. Molekul itu berputar pada
ikatan \ce{N-N} sampai pasangan-pasangan elektron bebasnya kira-kira tegak
lurus (konformasi gauche), tempat tumpang-tindihnya, dan tolakannya, kecil.
\end{solution}

\begin{solution}{pb:b2:frontier-orbitals:1}
\textbf{1.} Cincin beranggota lima dengan dua ikatan rangkap terkonjugasi
dan satu \ce{CH2}: cincin itu menahan satuan \omterm{def:b2:frontier-orbitals:diels-alder}{dienanya} dalam konformasi s-cis.
\textbf{2.} \omterm{def:b2:frontier-orbitals:diels-alder}{Dienofil}, melalui salah satu ikatan rangkapnya.
\textbf{3.} Kerangka norbornena yang terfusi dengan cincin siklopentena;
kedua ikatan $\sigma$ baru menghubungkan ujung-ujung \omterm{def:b2:frontier-orbitals:diels-alder}{diena} dengan kedua karbon
ikatan rangkap \omterm{def:b2:frontier-orbitals:diels-alder}{dienofil}.
\textbf{4.} \omterm{def:b2:frontier-orbitals:endo}{Adduk endo} (\omterm{def:b2:frontier-orbitals:endo}{aturan endo}): cincin \omterm{def:b2:frontier-orbitals:diels-alder}{dienofil} yang tersisa terletak
di bawah \omterm{def:b2:frontier-orbitals:diels-alder}{diena} dalam keadaan transisi.
\textbf{5.} Ya: reaksinya serempak, kedua ikatan terbentuk pada muka yang
sama dari masing-masing pasangan, dan geometri \omterm{def:b2:frontier-orbitals:diels-alder}{dienofil} dipertahankan.
\textbf{6.} Siklopentadiena sekaligus merupakan \omterm{def:b2:frontier-orbitals:diels-alder}{diena} yang baik (terkunci
s-cis) dan \omterm{def:b2:frontier-orbitals:diels-alder}{dienofil}, dan reaksinya serempak, tanpa zat antara bermuatan yang
perlu distabilkan.
\textbf{7.} Negatif: dua ikatan $\pi$ digantikan oleh dua ikatan $\sigma$.
\textbf{8.} Negatif: dua molekul menghasilkan satu.
\textbf{9.} $\Delta_r G^\circ = \Delta_r H^\circ - T\Delta_r S^\circ$ dengan kedua suku negatif: negatif
pada $T$ rendah, nol pada $T_i = \Delta_r H^\circ/\Delta_r S^\circ$, positif di atasnya.
\textbf{10.} Pada suhu kamar $T < T_i$: dimerisasi lebih disukai, walaupun
lambat. Di dekat titik didih dimer, kesetimbangan lebih menyukai monomer,
dan kesetimbangan itu tercapai dengan cepat.
\textbf{11.} Monomer meninggalkan campuran panas sebagai uap: mengeluarkan
sebuah produk membuat kesetimbangan terus bergeser ke arah monomer
(Le Chatelier).
\textbf{12.} Pada suhu kamar monomer itu berdimerisasi lagi; suhu dingin
memperlambat reaksi itu.
\textbf{13.} $0.62 - (-0.62) = 1.24|\beta|$.
\textbf{14.} $0.62 - (-0.35) = 0.97|\beta|$.
\textbf{15.} Siklopentadiena dengan anhidrida maleat: celah yang lebih kecil
memberikan stabilisasi keadaan transisi yang lebih besar, sehingga reaksinya
lebih cepat.
\textbf{16.} Keduanya terkonjugasi dengan ikatan \ce{C=C} dan menarik
elektron: seperti pada propenal, keduanya menurunkan tingkat $\pi^*$.
\textbf{17.} \omterm{def:b2:frontier-orbitals:frontier}{LUMO-nya}, yang mempunyai cuping pada karbon-karbon karbonil yang
berhadapan dengan karbon-karbon tengah \omterm{def:b2:frontier-orbitals:frontier}{HOMO} \omterm{def:b2:frontier-orbitals:diels-alder}{diena}.
\textbf{18.} Kerangka norbornena; cincin anhidrida di bawah ikatan rangkap
yang baru, berseberangan dengan jembatan \ce{CH2}.
\textbf{19.} Kedua karbon kepala jembatan dan kedua karbon yang membawa
anhidrida: empat pusat stereogenik.
\textbf{20.} Tidak: sebuah bidang cermin melalui jembatan \ce{CH2} dan di
antara kedua belahan molekul; adduk itu senyawa meso.
\textbf{21.} Keduanya cis pada \omterm{def:b2:frontier-orbitals:diels-alder}{dienofil}, dan adisi serempak mempertahankannya
tetap cis.
\textbf{22.} Kerangka yang sama dengan cincin anhidrida pada sisi jembatan
\ce{CH2}.
\textbf{23.} Tidak: mengubah endo menjadi exo berarti memutus dan membentuk
kembali ikatan \ce{C-C}, melalui reaksi balik, yang memerlukan pemanasan.
\textbf{24.} Air membuka anhidrida itu: asam dikarboksilat cis, dengan kedua
gugus \ce{COOH} endo.
\textbf{25.} \omterm{def:b2:frontier-orbitals:endo}{Adduk endo} mempunyai $\boldsymbol{4}$ pusat stereogenik.
\end{solution}
